\chapter{Discussão}
\label{cap:discussao}

Este capítulo interpreta os resultados apresentados nos Capítulos~\ref{cap:resultados} e~\ref{cap:validacao-tam} à luz da literatura revisada e do problema de pesquisa. A discussão organiza-se em torno de cinco eixos: a leitura dos resultados quantitativos da análise de sentimentos em Mossoró-RN, com destaque para o significado do padrão de neutralidade observado no ODS 9; a estabilidade dos padrões identificados quando a abordagem foi generalizada para outros municípios e outros objetivos; a triangulação entre a avaliação de aceitação tecnológica e os resultados quantitativos produzidos pelos dashboards; as limitações que delimitam o alcance das conclusões; e, por fim, as implicações dos achados para o conceito de Cidades Inteligentes e para a prática da gestão pública municipal. Antes de tratar dessas implicações, a metodologia é confrontada com os trabalhos relacionados, de modo a situar a contribuição da pesquisa no estado da arte.

\section{Discussão dos Resultados Quantitativos}
\label{sec:disc-quantitativos}

Os resultados obtidos para Mossoró-RN reforçam a premissa central desta pesquisa: avaliações públicas georreferenciadas podem ser utilizadas como fonte complementar de informação para a análise do desempenho dos ODS em nível municipal. Diferentemente de indicadores oficiais, que apresentam caráter agregado e atualização periódica limitada, os dados perceptivos analisados refletem experiências diretas dos cidadãos com serviços, infraestruturas e instituições específicas. Nesse sentido, a metodologia proposta não deve ser interpretada como um mecanismo de validação causal do desempenho dos ODS, mas como um instrumento exploratório e analítico, capaz de revelar padrões, gargalos recorrentes e sinais precoces de problemas urbanos e institucionais.

\subsection{O Padrão de Neutralidade no ODS 9: a Normalização de Problemas Crônicos}
\label{subsec:disc-neutros-ods9}

Um dos achados mais instigantes da análise quantitativa diz respeito à distribuição de sentimentos do ODS 9, marcada por uma elevada proporção de avaliações negativas (40,7\%) e, sobretudo, neutras (44,7\%), com apenas 14,6\% de avaliações positivas. A predominância simultânea de relatos negativos e neutros sugere um cenário de insatisfação estrutural associado à infraestrutura urbana e aos serviços de transporte. Contudo, é a elevada taxa de comentários neutros que merece interpretação mais cuidadosa, pois ela pode ser lida de duas formas complementares.

Por um lado, a neutralidade pode indicar relatos meramente descritivos, com baixo conteúdo avaliativo explícito -- comentários que registram a existência ou a localização de um equipamento sem emitir juízo de valor direto. Por outro lado, e de forma mais reveladora, a neutralidade pode refletir um processo de \textit{normalização de problemas crônicos}: deficiências recorrentes que, por sua persistência ao longo do tempo, deixam de ser percebidas como anomalias e passam a ser tratadas como parte esperada do cotidiano urbano. Sob essa leitura, o cidadão que descreve sem indignação um terminal rodoviário deteriorado não o faz por satisfação, mas porque a degradação tornou-se o estado de referência contra o qual a experiência é avaliada. A ausência de afeto explícito, nesse caso, é em si um sintoma da gravidade e da cronicidade do problema.

Essa interpretação tem consequências metodológicas importantes. Uma análise de sentimentos restrita à polaridade global classificaria a maior parte dessas avaliações como neutras e, consequentemente, subestimaria a extensão dos problemas de infraestrutura. É precisamente aqui que a opção por uma análise baseada em aspectos (\textit{aspect-based sentiment analysis}) demonstra seu valor: ao decompor cada avaliação nas sete dimensões do ODS 9 e ao identificar explicitamente menções problemáticas, o sistema recupera, de comentários aparentemente neutros, indícios de fragilidades persistentes que escapariam a uma classificação superficial. O fato de 56,9\% das análises contribuírem diretamente para o baixo desempenho -- proporção superior à dos comentários explicitamente negativos -- evidencia que parte substancial dos problemas foi extraída justamente de avaliações que a polaridade global não capturaria como negativas. Esse resultado corrobora a observação de \citeauthor{wankhade2022} \citeyear{wankhade2022} sobre os limites das abordagens baseadas exclusivamente em polaridade e reforça a pertinência de modelos generativos capazes de inferir aspectos implícitos no texto, conforme demonstrado por \citeauthor{agua2025} \citeyear{agua2025}.

\subsection{A Polarização no ODS 16: a Centralidade da Eficiência Institucional}
\label{subsec:disc-ods16}

O ODS 16 apresentou um padrão distinto, caracterizado por maior polarização entre avaliações positivas e negativas. Embora 32,8\% das análises tenham sido classificadas como positivas -- proporção mais que duas vezes superior à do ODS 9 --, a frequência e a recorrência de relatos negativos em dimensões críticas, especialmente instituições eficazes, acesso à justiça e segurança pública, foram suficientes para sustentar a pontuação muito baixa atribuída ao município. Esse contraste entre os dois objetivos é coerente com a natureza dos serviços avaliados: enquanto o ODS 9 se ancora em ativos físicos cuja degradação é contínua e cumulativa, o ODS 16 depende da qualidade do atendimento, da eficiência administrativa e da transparência dos processos -- elementos que tendem a gerar experiências mais episódicas e, portanto, avaliações mais emocionalmente carregadas, ora muito positivas, ora muito negativas.

A dominância da dimensão \textit{Instituições Eficazes e Responsáveis}, que concentrou 37,1\% das menções problemáticas, é o achado mais relevante para o desenho de políticas públicas. Ele indica que, em Mossoró, o principal fator explicativo do baixo ODS 16 não é a violência ou a impunidade em sentido estrito, mas a qualidade do serviço público -- atendimento demorado e burocrático, insuficiência de pessoal, piora percebida com a migração para canais digitais e deficiências de gestão. As correlações identificadas entre as dimensões sugerem a existência de efeitos sistêmicos: relatos de ineficiência institucional associam-se frequentemente a dificuldades de acesso à justiça, enquanto menções a violações de direitos tendem a coexistir com percepções de baixa transparência e de \textit{accountability}. Esses padrões reforçam a natureza transversal do ODS 16 e indicam que falhas institucionais não ocorrem de forma isolada, mas se propagam por diferentes esferas da governança pública. Tal leitura está alinhada à argumentação de \citeauthor{massey2022} \citeyear{massey2022} sobre a interdependência entre boa governança e o avanço da Agenda 2030.

\subsection{Padrões Transversais e Concentração Geográfica}
\label{subsec:disc-transversais}

A análise comparativa entre os dois objetivos revelou padrões estruturais comuns que transcendem suas especificidades: a falta de manutenção preventiva, a insuficiência de recursos humanos e financeiros, e a desconexão entre os serviços públicos e as necessidades cidadãs. A degradação por ausência de manutenção contínua afeta tanto a infraestrutura física (ODS 9) quanto os serviços institucionais (ODS 16), sugerindo que a fragilidade não é setorial, mas característica de um modelo de gestão reativo, em que a intervenção ocorre apenas após o colapso.

A dimensão geográfica acrescenta uma camada acionável a esses achados. A concentração de problemas em áreas centrais -- onde se localizam os órgãos públicos com maior volume de demanda --, em terminais de transporte e em regiões periféricas associadas à percepção de insegurança, evidencia que os problemas urbanos e institucionais possuem hotspots identificáveis. Essa característica é particularmente relevante para o ODS 9, no qual os problemas identificados estão fortemente vinculados a ativos físicos concretos, como o Terminal Rodoviário de Mossoró, repetidamente apontado como o principal foco de avaliações negativas. A possibilidade de associar percepções diretamente a locais específicos constitui uma vantagem da fonte georreferenciada e habilita a priorização espacial de intervenções, tema retomado na Seção~\ref{sec:disc-implicacoes}.

\section{Discussão da Generalização da Abordagem}
\label{sec:disc-generalizacao}

As evidências de generalização apresentadas nas Seções~\ref{sec:generalizacao-multiplos-ods} e~\ref{sec:generalizacao-geografica} permitem discutir até que ponto os padrões identificados em Mossoró são artefatos do caso específico ou manifestações de um comportamento estável da metodologia. A discussão a seguir distingue três aspectos: a estabilidade dos padrões em outros municípios, a sensibilidade da abordagem a diferentes níveis de desempenho e a fronteira entre generalização operacional e generalização substantiva.

\subsection{Estabilidade dos Padrões em Outros Municípios}
\label{subsec:disc-estabilidade}

A reaplicação dos modelos dos ODS 9 e 16 a quatro capitais do Nordeste, conduzida em caráter exploratório e com amostras reduzidas, produziu resultados qualitativamente consistentes com os de Mossoró. Em Fortaleza-CE, a execução do modelo do ODS 9 voltou a identificar a precariedade e a má gestão da infraestrutura de transporte e as falhas de mobilidade urbana como causas centrais do baixo desempenho, com 90\% das análises contribuindo para a pontuação muito baixa -- a mesma dimensão dominante observada em Mossoró. Em João Pessoa-PB, o modelo do ODS 16 reproduziu o padrão de centralidade institucional e de segurança pública, com 90\% das análises apontando problemas e predominância de sentimentos negativos. A repetição, em contextos urbanos distintos, da mesma estrutura de problemas -- infraestrutura de transporte para o ODS 9, instituições e segurança para o ODS 16 -- sugere que a abordagem captura regularidades que não são exclusivas de Mossoró, mas que se manifestam de forma análoga em municípios que compartilham desempenho oficial igualmente baixo.

Essa estabilidade qualitativa é relevante porque indica que as dimensões de análise definidas para cada ODS, os filtros semânticos e os \textit{prompts} especializados não foram excessivamente ajustados ao caso de Mossoró: eles transferem-se para novos contextos sem perda de coerência diagnóstica. Em termos da Design Science, trata-se de evidência de que o artefato possui um grau de generalidade que ultrapassa a instância que motivou seu desenvolvimento. Cabe sublinhar, ainda assim, que essas execuções em outras capitais permanecem exploratórias: por se apoiarem em poucas dezenas de avaliações, elas evidenciam a recorrência do padrão, mas não constituem diagnósticos validados de Fortaleza, João Pessoa, Natal ou Salvador. A única combinação de objetivo e município efetivamente analisada e validada nesta dissertação são os ODS 9 e 16 em Mossoró-RN.

\subsection{Sensibilidade a Diferentes Níveis de Desempenho}
\label{subsec:disc-sensibilidade}

Um achado mais sutil emerge da comparação entre municípios com diferentes pontuações oficiais. Nas execuções exploratórias dos ODS 3 e 4 -- conduzidas, como ressaltado, apenas para testar a generalização temática, e não como diagnósticos validados de saúde ou educação --, o comportamento da abordagem variou de forma informativa. Em Natal-RN, a análise do ODS 4 (classificado como ``Médio'') registrou distribuição de sentimentos mais equilibrada, com 44\% de avaliações positivas e 64\% das análises apontando problemas, refletindo uma realidade de contrastes em que coexistem focos de excelência e deficiências relevantes. Já em Fortaleza-CE, a análise do ODS 3 (também ``Médio'') flagou 90\% das análises como contributivas para o baixo desempenho, revelando uma insatisfação com os serviços de saúde mais intensa do que a pontuação oficial intermediária sugeriria.

Essa divergência merece interpretação. Ela indica que a métrica derivada das avaliações públicas -- o percentual de análises que apontam problemas -- não mantém correspondência linear com o índice oficial agregado. Tal descompasso não invalida a abordagem; ao contrário, ilustra que as duas fontes medem coisas distintas. O indicador oficial sintetiza múltiplas variáveis estruturais em uma pontuação, enquanto a análise perceptiva captura a experiência vivida com pontos de serviço específicos, que pode ser mais crítica do que a média sugere, sobretudo em equipamentos de alta circulação. Essa complementaridade -- e não a redundância -- é o que justifica o uso das avaliações como camada adicional de diagnóstico.

\subsection{Generalização Operacional e Generalização Substantiva}
\label{subsec:disc-operacional-substantiva}

É necessário, contudo, delimitar com precisão o que essas execuções demonstram. A generalização aqui evidenciada é de natureza técnica e operacional: ela comprova que o pipeline completo -- coleta na Google Places API, filtragem hierárquica, análise por aspecto, geração de JSON e produção de diagnóstico automatizado -- executa corretamente para diferentes objetivos e municípios, mediante apenas a substituição do arquivo de configuração, sem alteração do código-fonte. No caso dos ODS 3 e 4, essa substituição foi inclusive realizada por meio do formulário interativo de geração de configurações (Seção~\ref{sec:formulario-config}), o que reforça que a extensão da abordagem a novos objetivos está ao alcance de usuários sem conhecimento de programação. As execuções utilizaram amostras intencionalmente reduzidas, dimensionadas para confirmar a viabilidade do processo e controlar custos de API, e não constituem diagnósticos validados desses municípios nem conclusões sobre os domínios de saúde e educação.

Reitera-se, por isso, que apenas os ODS 9 e 16 em Mossoró-RN foram efetivamente validados e dispõem de dados suficientes para análise. A generalização substantiva dos achados -- a afirmação de que os mesmos problemas identificados em Mossoró ocorrem, na mesma magnitude, em outras cidades ou em outros objetivos -- exigiria estudos comparativos adicionais, com coleta em escala, validação local e análise das especificidades urbanas, institucionais e socioeconômicas de cada contexto. A distinção entre essas duas formas de generalização é essencial para a honestidade metodológica do trabalho: a pesquisa demonstra que a ferramenta é replicável e que os padrões qualitativos são estáveis nos casos exploratórios analisados, mas reserva a generalização empírica plena para investigações futuras, conforme detalhado no Capítulo~\ref{cap:conclusao}.

\section{Discussão da Aceitação Tecnológica}
\label{sec:disc-tam}

A avaliação de aceitação tecnológica conduzida com o instrumento TAM (Capítulo~\ref{cap:validacao-tam}) oferece uma perspectiva distinta e complementar à análise quantitativa: enquanto os Capítulos~\ref{cap:resultados} e a discussão precedente tratam da \textit{validade do conteúdo} produzido pelo sistema, a avaliação TAM trata da \textit{utilidade e usabilidade percebidas} da camada de apresentação por seus potenciais usuários. A articulação entre essas duas perspectivas constitui um exercício de triangulação metodológica.

\subsection{Triangulação entre Aceitação Percebida e Resultados Quantitativos}
\label{subsec:disc-triangulacao}

A triangulação revela uma convergência significativa. Os dashboards apresentam, do ponto de vista quantitativo, dados consistentes com os indicadores oficiais e diagnósticos que reproduzem fielmente os números agregados da coleta (Seção~\ref{sec:diagnosticos-automaticos}); do ponto de vista dos usuários, esses mesmos dashboards obtiveram Percepção de Utilidade média de 4,13 em escala de 5 pontos e satisfação de 4,6 tanto com a clareza dos dados quanto com a profundidade dos diagnósticos. Em outras palavras, a qualidade objetiva do conteúdo -- demonstrada pela consistência com o IDSC e pela fidelidade dos diagnósticos -- foi acompanhada por uma percepção subjetiva de utilidade igualmente elevada. Essa correspondência fortalece a confiança em ambos os resultados: a aceitação não se apoia em uma interface atraente que recobre conteúdo frágil, nem o conteúdo robusto é prejudicado por uma apresentação inacessível.

A triangulação também esclarece o papel relativo dos dois componentes da ferramenta. No ranqueamento de importância, os dashboards principais e os diagnósticos gerados pela IA foram considerados igualmente relevantes pela maioria dos avaliadores, e as opções mais selecionadas sobre a utilidade dos diagnósticos -- ``ajudam a identificar áreas críticas para intervenção'' e ``oferecem uma visão clara do progresso de Mossoró nos ODS'' -- correspondem exatamente às funções que a análise quantitativa demonstrou que o sistema cumpre: a identificação de hotspots (Seção~\ref{subsec:disc-transversais}) e a tradução de pontuações abstratas em causas concretas. Há, portanto, alinhamento entre o que o sistema efetivamente produz e o que os usuários percebem como seu valor.

\subsection{Convergência e Divergência entre Construtos}
\label{subsec:disc-construtos}

A leitura conjunta dos construtos do TAM com os achados qualitativos é igualmente esclarecedora. A Percepção de Facilidade de Uso (3,93) ficou ligeiramente abaixo da Percepção de Utilidade, com dispersão concentrada na curva de aprendizado inicial, e a única demanda técnica específica registrada -- a inclusão de filtros avançados, semelhantes aos de planilhas eletrônicas, na tela de análise detalhada -- converge com a única avaliação ``Regular'' atribuída à capacidade de \textit{drill-down}. Essa convergência entre uma resposta fechada (a nota ``Regular'') e uma resposta aberta (a sugestão de filtros) ilustra o ganho da triangulação interna ao próprio instrumento: a crítica não é um ruído isolado, mas um sinal consistente que aponta para um ponto específico e acionável de evolução da interface.

De modo análogo, a crítica de que os diagnósticos seriam ``superficiais e pouco acionáveis'', registrada por um avaliador, dialoga com a sugestão, feita por outro, de incluir ``propostas do que fazer para melhorar tal serviço''. Ambas apontam para a mesma lacuna -- a distância entre diagnóstico e prescrição -- e indicam que, embora o sistema já gere recomendações estratégicas, estas precisam ganhar maior proeminência e granularidade na interface. A convergência entre avaliadores de perfis distintos (gestão e academia) em torno dessa mesma demanda confere-lhe robustez e a qualifica como prioridade de desenvolvimento, e não como preferência idiossincrática. Vale ressaltar, ainda, que a presença dessas críticas em um conjunto majoritariamente favorável aumenta a credibilidade das avaliações positivas, ao indicar que os respondentes se sentiram à vontade para registrar discordâncias.

\section{Validação Metodológica à Luz da Literatura}
\label{sec:disc-validacao-metodologica}

\subsection{Eficácia da Filtragem Hierárquica}
\label{subsec:disc-filtros}

A estratégia de filtragem em múltiplas camadas mostrou-se determinante para a qualidade do corpus analisado. Os filtros geográficos eliminaram avaliações sobre estabelecimentos de municípios vizinhos capturadas pelo raio de busca, enquanto os filtros de relevância semântica garantiram o processamento apenas de comentários com conteúdo substantivo sobre os aspectos dos ODS. A exceção estratégica para locais de infraestrutura crítica -- que aceita todos os comentários independentemente de filtros de termos -- foi especialmente relevante: comentários altamente informativos sobre o Terminal Rodoviário, decisivos para o diagnóstico do ODS 9, teriam sido descartados por filtros mais restritivos. Essa engenharia de seleção de corpus é consistente com a observação de \citeauthor{wankhade2022} \citeyear{wankhade2022} de que o pré-processamento é fator crítico para a qualidade de análises de sentimentos, e equilibra precisão (relevância dos dados) e cobertura (quantidade de análises).

\subsection{Precisão da Análise com Modelos de Linguagem}
\label{subsec:disc-llm}

Os modelos de linguagem utilizados demonstraram capacidade consistente de identificar aspectos específicos dos ODS em textos livres, com scores de confiança médios superiores a 0,84, sem qualquer \textit{fine-tuning} para o domínio. Esse resultado alinha-se às conclusões de \citeauthor{larosa2025} \citeyear{larosa2025}, segundo as quais LLMs genéricos com \textit{prompts} bem elaborados podem rivalizar com modelos especializados em tarefas de análise de políticas de sustentabilidade, sobretudo quando a quantidade de dados rotulados é insuficiente para treinamento supervisionado. A abordagem de análise de sentimentos baseada em aspectos via LLMs, validada por \citeauthor{agua2025} \citeyear{agua2025} no contexto de avaliações de hotéis, confirma-se aplicável ao domínio de serviços públicos e infraestrutura urbana.

Cabe, porém, uma ressalva interpretativa: os scores de confiança fornecidos pelos modelos devem ser lidos com cautela, pois refletem a consistência interna das respostas geradas e não constituem, por si sós, uma validação externa da análise. A robustez dos resultados decorre, sobretudo, da agregação estatística de múltiplas análises e da recorrência de padrões ao longo de diferentes estabelecimentos e dimensões -- e não da confiança declarada em uma classificação individual.

\subsection{Análise Georreferenciada em Comparação com Redes Sociais}
\label{subsec:disc-georreferenciada}

A análise georreferenciada revelou vantagens importantes em relação a abordagens baseadas exclusivamente em redes sociais. Enquanto estudos anteriores, como os de \citeauthor{jaiswal2024} \citeyear{jaiswal2024} e \citeauthor{liu2023b} \citeyear{liu2023b}, capturam opiniões difusas e frequentemente desconectadas de experiências físicas verificáveis, o uso da Google Places API permitiu associar percepções diretamente a locais específicos -- terminais de transporte, órgãos públicos e equipamentos urbanos. Essa característica é particularmente valiosa para o ODS 9, cujos problemas estão ancorados em ativos físicos concretos, e para o ODS 16, no qual a concentração de avaliações negativas em poucos estabelecimentos institucionais reflete a centralidade desses órgãos na experiência cotidiana dos cidadãos, mesmo quando o número absoluto de locais é reduzido. A ancoragem espacial converte a análise de opinião em um diagnóstico territorializado, com endereço, requisito indispensável para a ação pública.

\section{Limitações do Estudo}
\label{sec:disc-limitacoes}

Os achados desta pesquisa devem ser interpretados à luz de limitações que delimitam seu alcance e que, simultaneamente, apontam caminhos para o aprimoramento futuro.

\subsection{Viés de Seleção e Exclusão Digital}
\label{subsec:disc-vies-selecao}

A limitação mais relevante é o viés de seleção inerente à fonte de dados. Avaliações online estão sujeitas à autosseleção dos avaliadores: usuários mais insatisfeitos ou mais engajados tendem a se manifestar com maior frequência, o que pode distorcer a distribuição de sentimentos. Mais grave, esse viés está intrinsecamente ligado à exclusão digital. Camadas da população com menor acesso à conectividade ou menor letramento digital -- frequentemente as mais vulneráveis e as mais dependentes de serviços públicos -- podem estar sub-representadas na amostra. Em consequência, a ausência de reclamações em determinadas áreas periféricas não deve ser interpretada automaticamente como satisfação, mas possivelmente como \textit{silêncio digital}. Essa constatação tem implicação prática direta: gestores que adotem a ferramenta devem combiná-la com métodos analógicos de escuta ativa, de modo a garantir equidade no monitoramento dos ODS e evitar que a alocação de recursos reproduza, em vez de corrigir, desigualdades de acesso.

\subsection{Cobertura Desigual e Temporalidade}
\label{subsec:disc-cobertura-temporalidade}

A distribuição desigual de avaliações entre estabelecimentos limita a cobertura em determinadas áreas urbanas, especialmente em equipamentos públicos menos digitalizados. Soma-se a isso a limitação de temporalidade: os comentários podem refletir situações passadas e não necessariamente o estado atual do serviço ou da infraestrutura. Como a plataforma não permite, na maioria dos casos, filtrar avaliações por período, a metodologia captura um \textit{snapshot} acumulado da percepção pública, e não uma série temporal. Essas limitações reforçam a necessidade de interpretar os resultados como indícios analíticos, e não como medições exatas, e motivam a proposta de coleta periódica discutida no Capítulo~\ref{cap:conclusao}.

\subsection{Viés de Confirmação dos Termos de Busca}
\label{subsec:disc-vies-termos}

Um aspecto metodológico adicional é o potencial viés de confirmação introduzido pela definição dos termos de busca. Por terem sido elaborados a partir de categorias tematicamente alinhadas aos ODS analisados, esses termos tendem a direcionar a amostra para estabelecimentos já contextualmente relacionados a essas dimensões, podendo reduzir a visibilidade de problemas em setores não contemplados e superestimar a prevalência de relatos negativos naqueles explicitamente buscados. Trata-se de uma tensão inerente à coleta orientada por objetivo: a mesma estratégia que confere foco temático ao corpus também restringe seu escopo. Estudos futuros poderiam mitigar esse efeito combinando termos mais amplos ou procedimentos de amostragem aleatória, de modo a contrabalançar a especificidade da busca dirigida.

\subsection{Limitações da Validação por Aceitação Tecnológica}
\label{subsec:disc-lim-tam}

A avaliação TAM, embora favorável, apoia-se em uma amostra de apenas cinco respondentes, recrutados por conveniência. Esse tamanho não permite inferência estatística sobre a população de potenciais usuários, e as médias reportadas constituem estatísticas descritivas sensíveis à variação de respostas individuais. Pela mesma razão, não foram calculados coeficientes de consistência interna nem testes de significância, cuja aplicação seria estatisticamente inadequada nesse contexto. Há, ainda, risco de viés de cortesia e de autosseleção, e a exposição dos participantes ao artefato limitou-se a uma única sessão, o que distingue a intenção de uso declarada do uso continuado real que o modelo se propõe a prever. A avaliação deve, portanto, ser lida como um estudo exploratório de aceitação, que fornece evidência inicial e estabelece um protocolo replicável, e não como uma validação psicométrica conclusiva -- conforme já delimitado na Seção~\ref{sec:ameacas-validade-tam}.

\subsection{Limitações Técnicas e de Custo}
\label{subsec:disc-lim-tecnicas}

Por fim, o uso extensivo da Google Places API implica custos crescentes com a escala da análise, e o \textit{rate limiting} imposto pela plataforma limita a velocidade de coleta, tornando análises de alta cobertura demoradas. O custo de \textit{tokens} dos modelos de linguagem, especialmente das variantes Pro, pode ser significativo em execuções com grandes volumes de comentários. Esses fatores não inviabilizam a abordagem -- que permanece substancialmente mais barata que pesquisas domiciliares --, mas condicionam o dimensionamento das coletas e justificam as amostras reduzidas adotadas nas execuções de generalização.

\section{Implicações para Smart Cities e Gestão Pública Municipal}
\label{sec:disc-implicacoes}

Para além de sua contribuição metodológica, os achados desta pesquisa possuem implicações concretas para o conceito de Cidades Inteligentes e para a prática da gestão pública municipal. Ao operacionalizar o uso de dados gerados pelos próprios cidadãos para a melhoria dos serviços urbanos, o trabalho alinha-se à missão central das Cidades Inteligentes: aumentar a eficiência governamental e a qualidade de vida urbana por meio de tecnologias digitais e participativas.

A primeira implicação é a viabilização de um diagnóstico municipal contínuo e de baixo custo. Os indicadores tradicionais dos ODS, embora rigorosos, são caros, temporalmente defasados e espacialmente grosseiros no nível municipal -- limitação especialmente crítica para objetivos vinculados à infraestrutura urbana e à efetividade institucional, nos quais a experiência cidadã constitui sinal valioso e tempestivo. A abordagem proposta complementa esses indicadores com uma camada perceptiva atualizável, capaz de sinalizar problemas antes que sejam captados por levantamentos oficiais. Esse potencial de monitoramento dinâmico é coerente com a visão de \citeauthor{fu2023} \citeyear{fu2023} sobre a colaboração entre humanos e modelos de linguagem na pesquisa urbana, na qual os LLMs atuam como instrumentos de ampliação da capacidade analítica do gestor público.

A segunda implicação é a priorização espacial e setorial de intervenções. A capacidade de territorializar problemas -- de apontar que o Terminal Rodoviário concentra avaliações negativas, ou que determinadas secretarias acumulam queixas de ineficiência -- permite que recursos escassos sejam direcionados aos pontos de maior impacto. No caso do ODS 9, a concentração de problemas em poucos equipamentos sugere que intervenções pontuais podem gerar efeito desproporcional sobre a percepção geral. No caso do ODS 16, a identificação da eficiência institucional como causa-raiz indica que medidas de melhoria do atendimento -- redução de filas, digitalização com suporte presencial adequado, capacitação de servidores -- podem produzir resultados mais imediatos do que programas de alto custo voltados a outras dimensões, em consonância com a relação entre boa governança e desenvolvimento sustentável discutida por \citeauthor{massey2022} \citeyear{massey2022}.

A terceira implicação diz respeito à transparência e ao controle social. Por serem compostos de arquivos estáticos e publicados gratuitamente na internet, os dashboards desenvolvidos democratizam o acesso a diagnósticos baseados em evidências, colocando-os ao alcance não apenas de gestores, mas também de pesquisadores, jornalistas e da sociedade civil. Essa abertura converte o diagnóstico em um bem público e cria condições para uma fiscalização cidadã mais informada do progresso municipal rumo às metas da Agenda 2030 \cite{onu2015}. A boa recepção da ferramenta entre profissionais da gestão pública, evidenciada na avaliação TAM, sugere que há demanda concreta por instrumentos dessa natureza.

Cabe reiterar, no entanto, a ressalva de equidade. Para que a adoção dessas tecnologias não aprofunde desigualdades, é indispensável que o diagnóstico digital seja combinado com mecanismos de escuta das populações sub-representadas na esfera online. Cidades verdadeiramente inteligentes não são apenas as que coletam mais dados, mas as que asseguram que esses dados representem o conjunto de seus cidadãos -- inclusive aqueles cujo silêncio, na ausência de cuidado metodológico, poderia ser confundido com satisfação.

\section{Comparação com Métodos Tradicionais}
\label{sec:disc-comparacao}

Em síntese, a metodologia proposta oferece vantagens claras sobre os métodos tradicionais de avaliação de ODS em quatro dimensões: velocidade (dias, em vez de meses), custo (substancialmente inferior ao de pesquisas domiciliares), granularidade (nível do estabelecimento, em vez de indicadores agregados) e especificidade (causas concretas, em vez de pontuações abstratas). Os métodos tradicionais, por outro lado, mantêm superioridade em representatividade estatística e em rigor de validação científica. Disso decorre que a abordagem aqui desenvolvida deve ser entendida como complementar, e não substitutiva, dos indicadores oficiais como o IDSC: seu papel é enriquecer o diagnóstico com a textura da experiência cidadã, e não reivindicar a precisão de uma medição populacional. É nessa complementaridade -- entre o rigor agregado do indicador oficial e a riqueza granular da percepção pública -- que reside a contribuição central desta pesquisa para o monitoramento dos ODS em nível municipal.
